\section{Introduction}

\subsection{Notations}
\begin{itemize}
	\item Observations are described by $p$ \hl{features} $\bm{X}=(X_{1},...,X_{p})\in\mathit{E}$ $ (E=\mathbb{R}^{p},...)$
	\item $\bm{X}_{i}=(X_{i1},...,X_{ip})$ is the features for observation $i$ $(1 \leq i \leq n)$
	\item $\mathit{Z}_i\in\{1,...,\mathit{K}\}$ is the \hl{group} of observation $i$
\end{itemize}

\subsection{The Mixture Model}
\textbf{Idea:} each group is described by its \textit{own} probability distribution.
\[
\bm{X} \mid \mathit{Z} = k \sim f(\mathbf{x},\theta_{k})=f_k(\mathbf{x})
\]

\textbf{How to read it:}
\begin{itemize}
\item $\bm{X}$ = feature vector describing one observation.
\item $Z=k$ = the observation belongs to group/cluster $k$.
\item $\bm{X}\mid Z=k$ = the features \emph{given that} the observation belongs to group $k$.
\item $\sim$ = ``is distributed according to''.
\item $f_k(\mathbf{x})$ = the probability distribution (density) for group $k$.
\item $\theta_k$ = parameters describing the distribution for group $k$.
\end{itemize}

\[
\fbox{
\parbox{0.8\textwidth}{
If an observation belongs to group \(k\), its features follow
the distribution \(f_k\) with parameters \(\theta_k\).
}}
\]

\textbf{Key idea:} Each group has its \emph{own} probability distribution:
\[
f_1(\mathbf{x}),;f_2(\mathbf{x}),;\ldots,;f_K(\mathbf{x}).
\]

\textbf{Important for clustering:} We observe $\bm{X}_i$, but $Z_i$ is typically
\emph{unknown}. We therefore use the observed features to infer which group
generated each observation.

\subsection{Mixing Proportions and Latent Group Membership}

The probability that an observation belongs to group $k$ is called the
\hl{mixing proportion}:

$$
p_k=P(Z=k),
\qquad k=1,\ldots,K.
$$

The mixing proportions satisfy

$$
p_k\geq 0
\qquad\text{and}\qquad
\sum_{k=1}^K p_k=1.
$$

Thus, $p_k$ represents the proportion (or prior probability) of observations
belonging to group $k$.

\medskip

\textbf{One-hot representation of the group membership:}

The group variable $Z$ can be represented by a vector
$\tilde{Z}=(\tilde{Z}_1,\ldots,\tilde{Z}_K)$ such that

$$
Z=k
\quad\Longleftrightarrow\quad
\tilde{Z}=(0,\ldots,0,1,0,\ldots,0),
$$

where the $1$ is in the $k$-th position. Therefore,

$$
\tilde{Z}_k=
\begin{cases}
1 & \text{if } Z=k,\\
0 & \text{otherwise.}
\end{cases}
$$

The membership vector follows a multinomial distribution with one trial:

$$
\tilde{Z}\sim\mathcal{M}(1,p_1,\ldots,p_K),
$$

and

$$
p_k=P(Z=k)=P(\tilde{Z}_k=1).
$$

\textbf{Key idea:} First, a group is selected according to the mixing
proportions $p_1,\ldots,p_K$. Then, the observation is generated from the
corresponding group distribution.

\subsection{Marginal Distribution of X: Mixture Density}
$\bm{X}$
Conditionally on belonging to group $k$,

$$
\bm{X}\mid Z=k\sim f_k(\mathbf{x}).
$$

Since $Z$ is generally unobserved, the marginal distribution of $\bm{X}$ is
obtained by combining all group distributions:

$$
\boxed{
f(\mathbf{x})
=
\sum_{k=1}^K p_k f_k(\mathbf{x})
}
$$

This is called the \hl{mixture density}.

$$
\boxed{
\text{Mixture density}
=
\sum
\text{mixing proportion}
\times
\text{group density}
}
$$

\subsection{Posterior Probability of Group Membership}

For an observed feature vector $\bm{X}=\mathbf{x}$, we want to compute the
probability that it belongs to group $k$:

$$
P(Z=k\mid \bm{X}=\mathbf{x}).
$$

Using Bayes' theorem:

$$
\boxed{
P(Z=k\mid \bm{X}=\mathbf{x})
=
\frac{p_kf_k(\mathbf{x})}
{\displaystyle\sum_{\ell=1}^K p_\ell f_\ell(\mathbf{x})}
=
\frac{p_kf_k(\mathbf{x})}{f(\mathbf{x})}
}
$$

\textbf{Interpretation:}
\begin{itemize}
\item $p_k$: how common group $k$ is (\textit{prior probability});
\item $f_k(\mathbf{x})$: how well group $k$ explains the observation
$\mathbf{x}$;
\item $P(Z=k\mid\bm{X}=\mathbf{x})$: the probability that the observation
belongs to group $k$ \emph{after observing} $\mathbf{x}$ (\textit{posterior
probability}).
\end{itemize}

\textbf{Memory aid:}

$$
\boxed{
\text{Prior probability}
\times
\text{How well the group explains }\mathbf{x}
\;\longrightarrow\;
\text{Posterior probability}
}
$$